\documentclass[14 pt,a4paper]{article}
\usepackage[a4paper, % размер страницы,
text={180mm, 260mm},% ширина и высота текста,
left=15mm, top=15mm]{geometry} % размеры пустых полей,
\usepackage[utf8]{inputenc} %кодировка текста,
\usepackage{cmap}  %кодировка pdf-документа,
\usepackage[english,russian]{babel} % поддержка русского языка,
\usepackage{indentfirst} % корректировка отступов,
\usepackage{amssymb} % математические символы,
\usepackage{amsmath} % математические конструкции,
\usepackage{graphicx}

\usepackage[backend=biber,style=gost-numeric]{biblatex}
\usepackage{xcolor}

\begin{document}
	
\section{Титульный лист}
\section{Актуальность} 
За последние годы наблюдается несколько существенных факторов роста интереса к задачам автоматического обнаружения аномалий:
\begin{itemize}
	\item во-первых, растёт спрос на автоматизацию в самых разных отраслях — от производства до инфраструктуры;
	\item во-вторых, усиливаются требования к аналитике в режиме реального времени: предприятия хотят получать информацию мгновенно и действовать на опережение;
	\item в-третьих, достижения в области ИИ и машинного обучения сделали возможным применение сложных моделей даже в критичных промышленных задачах.
\end{itemize}
Однако при внедрении есть и серьёзные проблемы, которые нельзя игнорировать:
\begin{itemize}
	\item высокие капитальные и эксплуатационные затраты современных систем компьютерного зрения;
	\item вопросы конфиденциальности и безопасности данных — особенно когда видеопотоки и изображения содержат чувствительную информацию.
\end{itemize}
Эти факторы определяют постановку задачи: требуется решение, которое эффективно, экономично и безопасно интегрируется в реальные процессы.
\section{Отрасли применения}
Применение технологии широкое — перечислю ключевые области и поясню один-два примера для каждой:
\begin{itemize}
	\item Металлургическое производство — автоматический контроль поверхностей: обнаружение трещин, вкраплений, дефектов проката, что повышает качество и снижает брак.
	\item Железнодорожный транспорт — мониторинг рельсов и подвижного состава для своевременного выявления износа и повреждений.
	\item Диагностика линий электропередач — раннее обнаружение коррозии, повреждений изоляции и посторонних предметов, что повышает безопасность сети.
	\item Гражданское строительство — визуальный контроль состояния конструкций, фасадов, мостов и инженерных сооружений для прогнозирования неисправностей.
	\item Агропромышленный сектор — сортировка продукции и выявление брака, что ускоряет обработку урожая и снижает потери.
\end{itemize}
\section{Цели и задачи}
...

\section{Описание предлагаемого решения}
Подход базируется на методах глубокого обучения, нацеленных на обучение по нормальным (базовым) данным без необходимости разметки дефектов. Главные черты решения:
\begin{itemize}
	\item Использование моделей реконструкции или представлений. Система обучается восстанавливать изображение или его признаки. Участки, которые плохо восстанавливаются, трактуются как аномальные.

	\item Обучение на распределениях нормальных данных. Модель формирует математическое представление того, как выглядит объект в исправном состоянии, и сравнивает новые изображения с этим распределением.

	\item Извлечение многоуровневых признаков. Глубокие сверточные сети позволяют выделять как глобальные формы, так и мелкие текстурные элементы, что важно для обнаружения дефектов.

	\item Использование современных архитектур. Это могут быть автоэнкодеры, вариационные автоэнкодеры, нормализующие потоки, трансформеры или модели с встраиваемыми картами внимания.

	\item Работа с небольшим количеством аномалий. Так как реальные дефекты встречаются редко, система концентрируется на изучении норм и выявлении отклонений, не требуя обширной разметки.

	\item Гибкость и адаптивность. Модель может быть переобучена под новый объект или тип поверхности без полного изменения архитектуры.
\end{itemize}
Такая комбинация дает сбалансированный результат: высокая чувствительность к аномалиям при приемлемых требованиях к разметке.

Универсальная модель (детали реализации)
Ключевые блоки и их назначение:
\begin{itemize}
	\item Автоматическая подготовка данных — сбор и фильтрация нормальных образцов, аугментация для повышения вариативности;
	\item Обучение на конкретных объектах — каждая модель адаптируется к типичной внешности изделия, что улучшает точность;
	\item Обучение на неразмеченных данных — используются подходы типа реконструкции, плотностного моделирования или контрастивного обучения;
	\item Визуализация найденных дефектов — тепловые карты, маски и bounding-box’ы для удобства оператора;
	\item Интерфейс оператора — дашборд с возможностью фильтрации, метрик и ручной разметки для последующего дообучения;
	\item Расширяемый модульный подход — независимые модули предобработки, модели детекции и интерфейса упрощают поддержку и интеграцию.	
\end{itemize}

\section{Преимущества и недостатки глубокого обучения}
Преимущества:
\begin{itemize}
	\item Автоматическое извлечение признаков — не требуется ручная разработка детекторов;
	\item Высокая точность на сложных текстурах и неоднородных поверхностях;
	\item Устойчивость к шуму и вариациям в освещении и ракурсе при корректной подготовке данных;
	\item Эффективность в условиях неразмеченности — система обучается на нормальных данных и выявляет отклонения;
	\item Масштабируемость — модель можно перенастроить под новые объекты без полной переработки.
\end{itemize}
	
Недостатки:
\begin{itemize}
	\item Требовательность к данным. Даже без примеров дефектов нужен большой объем «нормальных» образцов, отражающих вариативность процесса. Решение: организовать непрерывный сбор данных и использовать аугментации.

	\item Вычислительная сложность. Обучение требует ресурсов. Решение: использовать гибрид — обучение на сервере и оптимизированная инференс-версия для edge.

	\item Низкая интерпретируемость. Модель может давать сигнал без очевидного объяснения. Решение: внедрение модулей визуализации и инструментов объяснимости (например, attention-карт).

	\item Сложность разработки. Длительные эксперименты и настройка гиперпараметров. Решение: стандартизировать пайплайн и автоматизировать поиск гиперпараметров.
	
	\item Проблемы переобучения. Особенно на узких выборках. Решение: контролируемое тестирование на независимых и симулированных аномалиях, регуляризация и мониторинг производительности в продакшне.
\end{itemize}

\section{Флагманы отрасли}
\begin{itemize}
	\item NVIDIA. Один из мировых лидеров по аппаратному обеспечению для deep learning и computer vision — её GPU и платформы ускоряют обучение и инференс CV-моделей. 
	\item Intel. Предлагает решения для CV на уровне железа и софта: процессоры, акселераторы и инструменты (например, SDK) для edge- и embedded-приложений. 
	\item Microsoft. Через облачные сервисы и инструменты AI помогает компаниям внедрять CV (распознавание изображений, аналитика, автоматизация) без необходимости создавать собственную инфраструктуру. 
	\item Google. Имеет сильные AI- и ML-команды, развивает решения для обработки изображений, видео, и других визуальных данных — применимо в самых разных областях. 
	\item Cognex. Серьёзный игрок в индустриальном/промышленном компьютерном зрении: системы машинного зрения, визуального контроля, автоматизации на производстве и логистике. 
	\item Basler AG	Производитель “умных” камер и систем визуального контроля, востребованных для автоматизации, инспекции, роботов и промышленности. 
	\item Megvii (и платформа Face++). Один из крупнейших провайдеров CV-алгоритмов (чёрно-белое распознавание лиц, аналитика) — особенно заметен на рынках Китая и Азии
\end{itemize}

\section{Список литературы}.

\end{document}